% Transcription — fonds Alexandre Grothendieck, shelfmark 29, batch 1 (pages 1-20).
% Established from the facsimile archives/batches/29/batch-01.pdf (a mirror of
% https://grothendieck.umontpellier.fr/29.pdf), per the skill
% transcribe-grothendieck.
%
% Page 6 is the blank back of the carbon of page 7 — absent.
%
% The batch holds two pieces, and they belong together. Pages 1-11 are the
% 1967 correspondence with J.-P. Murre over the exposé on the tame fundamental
% group that Murre was writing up for SGA 1: Murre's letter of 29 March 1967
% sending in the manuscript with his own list of misgivings (pages 4-5),
% Grothendieck's reply listing the misstatements he blames on his own sketchy
% notes (pages 7-8), his letter of 29 April 1967 enclosing a fresh sketch on
% ramification data and proposing a plan in six sections for the talk (pages
% 9-11), and Murre's letter of 16 May 1967 saying he has not grasped the new
% set-up (pages 2-3). The archive's order is not the chronological one.
%
% Pages 12-20 are that enclosed sketch: a typescript headed "1. Données de
% ramification sur un schéma", defining a ramification datum as a pair
% R = (C, s) — a fibred category over the étale site with a cartesian functor
% to relative schemes — then Rev^R(S), the geometric realisation functor r, and
% three examples of increasing generality, the third of which is retyped: page
% 18 states it through a factorisation of s, page 19 replaces that by a
% condition on the automorphism groups. Both moutures are here, and they are
% not merged. Page 20 opens the Kummer cases.
%
% Page 17 is the back of one of those sheets and carries two inked arrows
% belonging to the example on the recto — transcribed, briefly.
%
% The letters are in English, the typescript in French. Both stay as written.
%
% The dating is the inventory's own, for the whole shelfmark. The three letters
% here date themselves, and 1967 falls inside it.
%
% Pass: Opus 5 (claude-opus-5), 2026-08-22 — first pass, unchecked against the pages by a human.
\documentclass[11pt,a4paper]{article}
% Le préambule commun à toutes les transcriptions du fonds Grothendieck.
%
% Il tient en une idée : **l'incertitude se compose, elle ne se résout pas.**
% Une transcription qui lisse un mot douteux a détruit la seule chose pour
% laquelle on la faisait — un lecteur doit pouvoir savoir, à chaque endroit,
% ce qui était sur le feuillet et ce qui a été ajouté. Les six macros qui
% suivent sont donc l'appareil critique, et non de la décoration.
%
% Les mêmes commandes sont reconnues par `scripts/render.mjs`, qui produit la
% vue de lecture du site. Une macro ajoutée ici doit l'être aussi là-bas,
% faute de quoi le rendu échouera bruyamment — ce qui est voulu : un rendu
% silencieusement faux est pire qu'un rendu absent.

\ProvidesPackage{grothendieck}[2026/08/08 Transcriptions du fonds Grothendieck]

\usepackage{amsmath,amssymb,amsthm}
\usepackage{mathrsfs}
\usepackage{stmaryrd}
% Les diagrammes commutatifs sont partout dans ce fonds : c'est la langue
% même de la géométrie algébrique de Grothendieck, et une transcription qui
% les rendrait en prose ne transcrirait rien.
\usepackage{tikz-cd}
% Français par défaut — c'est la langue des feuillets — mais l'anglais est
% chargé aussi : la traduction et le résumé sont des documents anglais, et la
% typographie française (espaces avant les deux-points) y serait une faute.
% Ces documents basculent par \selectlanguage{english} en tête de corps.
\usepackage[english,french]{babel}
\usepackage{fontspec}
\usepackage{geometry}
\geometry{a4paper,margin=25mm,marginparwidth=28mm}
\usepackage{marginnote}
\usepackage{xcolor}
% ulem, not soul. `soul`'s \st and \ul are built on a character-by-character
% scan that XeLaTeX's font machinery breaks: the document fails with an
% undefined control sequence rather than degrading. ulem does the same two jobs
% and is Unicode-engine safe. `normalem` stops it hijacking \emph, which the
% transcriptions use for Grothendieck's own emphasis.
\usepackage[normalem]{ulem}
\usepackage{hyperref}
\hypersetup{colorlinks=true,linkcolor=black,urlcolor=black,pdfborder={0 0 0}}

% --- Métadonnées du lot -----------------------------------------------------
% Reprises telles quelles par le script de rendu, qui les affiche en tête de
% la vue de lecture. Elles fixent ce que le fichier prétend couvrir : sans
% elles, une transcription flotte sans point d'attache dans un fonds de seize
% mille feuillets.

\newcommand{\folder}[1]{\gdef\grothFolder{#1}}
\newcommand{\batch}[1]{\gdef\grothBatch{#1}}
\newcommand{\leaves}[2]{\gdef\grothFirst{#1}\gdef\grothLast{#2}}
\newcommand{\foldertitle}[1]{\gdef\grothFoldertitle{#1}}
\gdef\grothFolder{?}\gdef\grothBatch{?}\gdef\grothFirst{?}\gdef\grothLast{?}\gdef\grothFoldertitle{}

% --- Le résumé --------------------------------------------------------------
%
% Il ouvre la lecture modernisée, sous le titre « Résumé », et il est composé
% en petit. Ce n'est pas un ornement : c'est le seul passage du document qui
% ne soit pas des mathématiques, et un lecteur doit pouvoir le reconnaître
% avant de l'avoir lu — puis le sauter s'il connaît déjà le sujet, ou s'y
% arrêter s'il ne le connaît pas. Le filet qui le suit marque la reprise du
% corps.

\newenvironment{resume}
  {\small\setlength{\parskip}{0.45em}}
  {\par\medskip\hrule\medskip}

% --- Mots-clés ---------------------------------------------------------------
%
% En anglais, en clôture du résumé de la lecture modernisée : le vocabulaire
% moderne sous lequel chercher ce que les feuillets construisent. Le manifeste
% du site (`npm run manifest`) les extrait pour étiqueter la cote entière —
% cette ligne est la source unique des tags, il n'y a pas de fichier à côté.
\newcommand{\keywords}[1]{\par\smallskip\noindent
  {\footnotesize\textsc{Keywords} --- \textit{#1}}\par}

% --- L'appareil critique ----------------------------------------------------

% Le numéro de feuillet, en marge. C'est l'ancre qui permet de régler un
% désaccord sur une formule en regardant *un* feuillet plutôt qu'une chemise
% de six cents. Le site s'en sert aussi pour tourner les pages du fac-similé
% quand on fait défiler la transcription.
\newcommand{\leaf}[1]{%
  \marginnote{\footnotesize\color{gray!70}\textsf{#1}}%
  \label{leaf:#1}\ignorespaces}

% Illisible : on ne devine pas. Un mot inventé qui se lit comme les autres est
% le pire résultat possible.
\newcommand{\ill}{\textcolor{red!60!black}{[\,\ldots\,]}}

% Lecture douteuse : le mot est donné, et signalé comme incertain.
\newcommand{\uncertain}[1]{\uline{#1}}

% Ajout de l'éditeur — une lettre manquante, un mot sous-entendu.
\newcommand{\add}[1]{\textcolor{blue!50!black}{[#1]}}

% Biffé par Grothendieck lui-même. Ce qu'il a rayé fait partie du document :
% une rature dit souvent où la pensée a changé de direction.
\newcommand{\struck}[1]{\sout{#1}}

% Note du transcripteur, sur l'état matériel du feuillet.
\newcommand{\note}[1]{\par\smallskip{\small\color{gray!60!black}\textit{[#1]}}\par\smallskip}

% Note marginale de Grothendieck, distincte du corps du texte.
\newcommand{\marginal}[1]{\par\smallskip\noindent\hspace{1em}%
  {\small\color{gray!60!black}#1}\par\smallskip}

% --- Titre ------------------------------------------------------------------

\AtBeginDocument{%
  \begin{center}
    {\large\bfseries Fonds Alexandre Grothendieck --- cote n\textsuperscript{o}~\grothFolder}\\[2pt]
    {\itshape\grothFoldertitle}\\[4pt]
    {\small feuillets \grothFirst--\grothLast\ (lot \grothBatch)}\\[2pt]
    {\footnotesize\color{gray!60!black}Transcription établie d'après le fac-similé numérisé
     par l'Université de Montpellier.\\
     Les lectures incertaines sont soulignées, les illisibles notées [\,\ldots\,],
     les ajouts entre crochets.}
  \end{center}
  \vspace{1em}}

\endinput


\folder{29}
\batch{1}
\pages{1}{20}
\dating{1959-1969}
\watermark{Édition de démonstration}
\foldertitle{Groupe fondamental [Autour de SGA 4 et SGA 7] : lettres (1959, 1967, 1969), tapuscrits et copies de tapuscrit annotés (s.d.), notes manuscrites (s.d.)}

\begin{document}

\section*{Correspondance avec J.-P. Murre sur le groupe fondamental modéré (pages 1 à 11)}

\page{1}
Titre porté à l'encre en haut de la chemise, souligné : « Tame ramification ».
\note{c'est le titre de la main de Grothendieck ; celui de l'inventaire, entre
crochets, est des archivistes}

\subsection*{J.-P. Murre à A. Grothendieck, Voorschoten, 16 mai 1967 (pages 2 et 3)}

\page{2}
Voorschoten, May 16 1967. Palestrinalaan 11.

Dear Grothendieck,

I am trying to work out your letter with indications. However I have many
difficulties and my progress is very slow. Usually I could find out from your
manuscripts for what purpose you developped things in such or such a way, but
this time I doubt whether I have grasped the essential points. I can understand
your set-up only against the lemma of Abhyankar as background, namely instead
of giving directly the covering $r(X)$ (in your notation) of $S$, you give the
``behaviour'' of the covering over other coverings, namely the $\underline{s}(M)$,
which are of a particular type; apparantly the behaviour over all these
$\underline{s}(M)$ together gives better information than the $r(X)$ alone. Is
this intuitively a reasonable way of thinking over your new set-up?

My main difficulty is that I don't see that for a Galois covering the extra
structure is uniquely determined. First the concept of Galois covering itself:
After the abstract definition and some variations you write: ( I assume that
you have a copy) \emph{On peut enfin l'expliciter en revenant à la définition
de $\mathrm{Rev}^{\underline{R}}(S)$, et en exigeant que les revêtements des
$\underline{s}(M)$ correspondants à $X$ soient principaux de groupe $G$.} Now
in all your examples there is also a group working on the $\underline{s}(M)$,
for instance in example 1 a group $\mathcal{G}$ working on $Z$, and this group
acts also on the Galois covering $X$. It seems to me that one must require that
the actions of $G$ and $\mathcal{G}$ commute; is this correct? On the same and
the next page of your letter you explain (for example 1) that the functor
$\underline{r}$ (geometric realization) is fully faithful (restricted to Galois
coverings) under certain conditions; this I don't understand at all ( and
\struck{makes} this makes me feel very uncertain, since I expect that this is
an important point if one wants to see what is going on).
\marginal{corresponding with}
\note{« corresponding with », à l'encre bleue en marge, vise « the Galois
covering $X$ » ; les lettres grecques et scriptes de la lettre sont de même
ajoutées à la main, la machine ne les portant pas}

Some other points: in the final formulation of example 3 you introduce a sheme
of groups $G(M)$; \struck{with this formulation} do the objects in example 3a,
with this formulation, consist out of \emph{couples} $M=(\underline{a},
\underline{n})$ and $G(M) = \mu_{\underline{n}}$? Also you write:
\emph{$\mathrm{Hom}_{S'}(M',M)$ est vide sauf si $\underline{n}' \geqslant
\underline{n}$, auquel cas\dots{}etc\dots}, it seems to me that this should be
\dots\emph{sauf si $\underline{n}' = \underline{m}\cdot\underline{n}$, auquel
cas\dots} and where $\underline{m} = (m_i)_{i \in I}$ is a set of positive
integers.

\page{3}
I realize that I cause you a lot of trouble and work by asking every time for
further explanation and that my assistance in writing the exposé is not very
large. Therefore I suggest that you answer this letter only in case it is clear
from my questions that I have really a misconception of your new set-up (and it
would not surprise me if I have\dots). Otherwise I shall try to work out your
theory as good as I can and make a preliminary manuscript which we can discuss
when I am at Bures (I plan to come June 2). In any case, I hope to complete the
definite version during my stay in Bures.

With kind regards, Sincerely yours, J. P. Murre.

\subsection*{J.-P. Murre à A. Grothendieck, Voorschoten, 29 mars 1967 (pages 4 et 5)}

\page{4}
Voorschoten, March 29 1967. Palestrinalaan 11.

Dear Grothendieck,

Enclosed I send you --- finally ---, in two different covers, the manuscript on
the tame fund.\ group together with the papers which I borrowed from you.

As to the terminology. I have used \emph{scheme} instead of \emph{preschema}
(without mentioning it, because I expect this will be a general remark in
SGA I) but I have not used ``doublet'' because I don't know your final decision
on this point.
\marginal{stroke}

References. I have made freely references to EGA (also to EGA IV) but
\emph{only to} \emph{SGA 1960/61} (with one notable misleading : a reference to
SGA X, Cor 3.3 etc, the purity theorem, is a misleading because the proof is
only in SGA 1962). I don't know the ``new numbering'' of SGA (you mentioned
something about it, but I don't remember) so I have used still SGA I, SGA V etc
for the 1960/61 séminaire.
\marginal{SGA 1}

I must admit that there are still many points which are not satisfactory in the
exposition. Below are some comments, motivations and questions:

p.3 Hyp.1 --- Perhaps it would have been nicer to avoid a little longer this
assumption. However there is in SGA 1960/61 only a statement about quotients
and existence of quotients in case the group is an \emph{ordinary} finite group
(it is only in SGA D that the theorems from TDTE III are proved and the proof
is long and subtle!). This is the reason for making the assumption Hyp.1
already here.

p.5 line 6 below: the emphasis on $H$ ordinary group : same reason.

p.6 (Cor 1.11) \emph{In fact this is a general question} : I have hesitated a
little bit with the word ``covering'' for ``revêtement'' (of course: already
used in the title!) Originally revêtement was used as translation \struck{for
covering} for covering (in Lang's papers)! But now there may be confusion with
``covering in some topology'' for instance ``covering in the étale topology''
which is something entirely different. What do you think about this? I don't
see what
\marginal{O.K.}

\page{5}
word to use else.

p.19 etc : I expected lemma 4.3 to be almost trivial but when I tried to write
a proof I did not see an easier one. Also: the proof of lemma 4.4 is kind of
subtle ( but this could be expected).

p.26. line 1: I shall try to find a precise reference.

line 3: the condition : \emph{irreducible in Zariski topology} is unpleasant
but I have not been able to avoid it.

last line in 3) : it seems necessary to mention also the points $s'$.

p.33, 34 etc.\ --- I did not see an easier proof for prop 7.4.

p.43 theor.\ 9.4 --- You have mentioned that the theorem in this form is not
satisfactory. For instance the trivial case with constant fibers and arbitrary
crossings escapes. However I don't see how to get a satisfactory statement.

p.45. line 5 etc.\ --- Is this indication sufficient?

I have not written an appendix on the Lefschetz theory. If you think this is
useful then I can still try.

It is not necessary to return the manuscript, I have a carbon copy so you can
refer to pages and numbers in your comments. Please give me, if necessary,
indications for the way the final manuscript must be typed ( is the book a
\struck{foto} photo copy of the manuscript?) Also give me a deadline when you
must have the manuscript because I am afraid that I may have already delayed
the publication date of your first volume.

With best regards, Sincerely yours, J. P. Murre.

\subsection*{A. Grothendieck à J.-P. Murre (pages 7 et 8)}

\page{7}
Dear Murre,

Thanks you very much for your notes on the tame fundamental group, which I at
last finished reading. I see you wrote them with much care, and I am all the
more sorry that by my own fault, there is a number of misstatements which, I am
affraid, will force you to do a serious recasting of the whole exposition. My
notes definitely where too sketchy, and my oral explanations, I am affraid,
partly wrong, which induced you into error a few times. Here the most serious
drawbacks.

1.16.\ is false already when $H$ is the unit group and when there is a single
$a$, say $a = b^{n}$, $Y' = Y\,T/(T^{n}-a)$. Then $Y = S$, and a morphism of
$Y$ into $Y'$ compatible with $H = e \to H'$ is just a section of $Y'$, which
exists indeed; however $H \to H'$ is not surjective. 3.6.\ is equally false, as
you see by the previous example, using the given section to define an
$H'$-morphism $H' \to Y'$ which is not an isomorphism. As a consequence, the
proof in your notes of 3.7.\ breaks down (as it uses 1.16) and so does the
proof of 3.8. (I did not try to check 3.8.\ by some different proof).
\note{$Y\,T/(T^{n}-a)$ : lire $Y[T]/(T^{n}-a)$, la machine ne porte pas de
crochets ; de même les flèches sont laissées en blanc dans tout le lot et
rétablies ici}

I am affraid \struck{the proof of} 6.4.\ is false as stated, and that the
statement is correct only if the $D_i$ are regular. Indeed, the end of the
proof seemed to me very dubious; be careful that the intertia groups are
determined only up to interior automorphism! There is however a
\emph{(tautological)} generalization of the theorem for regular $D_i$,
corresponding to the data of a single divisor $D$ with normal crossings, and a
variant of the notion of tame ramification for such a divisor, by demanding
that the coverings should be tamely ramified locally for the étale topology for
the family of local irreducible components of $D$; it is this notion of
tameness which would seem more adapted to the situation of par.\ 9.

\page{8}
The proof of 7.1.\ is not correct, when you conted on line $-9$: there remains
to be proven the following: \dots{} Alraedy when $D=0$, the proof here would
have to introduce connected étale coverings which are \emph{not Galois}! This
very strongly suggests that a notion of tame ramification should be introduced
also \struck{in} for non Galois coverings. The same remark applies to the proof
of 10.1. Maybe you could get along some way in 7.1.\ using the normality
assumptions, but I am convienced that these assumptions are anyhow artificial,
as well as the assumption that the $D_i$ should be reduced somewhere. You do
not seem to make any use of these facts, really.

Also, one feels that 7.5.\ should be generalized to the case of a tamely
ramified covering, and that it should come out trivially once the generalities
have been dealt with properly.

I hope that the theory will come out more clearly and correctly by devoting
some care to generalities on ramification data (not necessarily of Kummer
type). I will try to write something up within the next days. Please excuse me
for the trouble I caused you by not learning my lesson well enough before I put
you to work!

It would be very nice indeed to have an appendix on Lefschetz theorem for the
tame fundamental group, and it should not be hard to write it. However, if
would be safer to wait till the general theory of tame ramification is written
up!

Sincerely your's

\subsection*{A. Grothendieck à J.-P. Murre, 29 avril 1967 (pages 9 à 11)}

\page{9}
\uncertain{29.4.67}

Dear Murre,

I am sending you enclosed the sketch which I promised on ramification data, as
well as your manuscript together with my annotations. As the paragraphs 1 to 5
will have to be changed entirely, there is maybe not much point for you to look
up to carefully my annotations on these. I would only suggest that you put more
formulas, arrows etc on \struck{separate} separate lines, for easier reading;
that you normalize the standard references to SGA, EGA (for instance, no need
to indicate SGA when reference is being made to the seminar 1960--61 in which
your talk is going to be included itself, reference as (XI 6.1) is then
sufficient); that you be cautious when using a locution like ``if primitive
$n$.th roots exist over $S$'', which does not mean much by itself, unless you
make clear what you really mean. I would propose the following tentative plan
for your talk:

\begin{enumerate}
  \item Kummer coverings with respect to a given set of sections.
  \item Ramification data.
  \item Properties of the catégory of coverings of type $\underline{R}$. The
    group $\pi_1^{\underline{R}}(S,\;)$.
  \item Functorial behaviours.
  \item Galois coverings of type $\underline{R}$.
  \item Abhyankars theorem (for divisors with normal crossings).
\end{enumerate}

etc

The sections 6 to 10 essentially the same as in your notes, just taking into
account my observations, and changing inner references according to the changes
made in 1.\ to 5. A paragraph 11 could be an appendix on the tame variant of
Lefschetz's theorem. --- Par.1 is essentially your previous par.1; you should
check exactly what is needed in it for what
\note{le troisième point porte $\pi_1$ sous la forme $\genfrac{}{}{0pt}{}{R}{1}(S,\;)$,
l'exposant $\underline{R}$ étant frappé au-dessus de l'indice 1 faute de pouvoir
le placer autrement}

\page{10}
is to follow. Moreover, it would be reasonable to include the fact that when
$S$ is normal and the $V(a_i)$ reduced, then the $Z^{\underline{n}}_{\underline{a}}$
are normal (which now is included in your par.5), so that you can refer to this
fact when mentioning examples in \struck{sections} 2.\ to 5. The paragraphs 2
to 5 in this plan correspond approximately to the sections 1.\ to 4.\ in my
sketchy notes. It will make some work, I am affraid, to put it into shape, but
I have the feeling the notions developped with some generality can be useful
still in other contexts than just for the kummer type ramification. If you
think it would be much simpler to write these sections by restricting
throughout to kummer type ramification, please feel free to do so, (and refer
to the possible generalization by a suitable remark). I doubt however it will
be much simpler, as there are at least two useful variants of the notion
besides the one of your notes (cf ``particular case b'' in my sketch).

Some more specific observations. In theorem 6.4.\ (which should be formulated,
I believe, in terms of tame ramification in the large sense, cf my notes),
there is no need to look at the residue characteristics of the non maximal
points of the $D_i$. It will turn out a posteriori (and this should be pointed
out in a reamrk or corollary) that if for every maximal point $s$ of every
$D_i$, the inertia group at $s$ is prime to the charactéristic of $k(s)$, then
the order of this group must be automatically prime to the residue
characteristic at any specialization of $s$. (This statement is non empty only
when $s$ is of char.\ zero, of course). In fact, the proof in your notes (after
reduction to the strict local case, which is immediate) works out unchanged,
provided you observe that if $A$ is a regular local ring, $(x_i)_{1 \leqslant i
\leqslant r}$ part of a regular system of parameters, and $(n_i)_{1 \leqslant i
\leqslant r}$ a system of integers $\geqslant 1$, not noecessarily

\page{11}
prime to the residue characteristic, then the ring
\[
  B \;=\; A\,[T_1,\dots,T_r]\,/\,(T_1^{\,n_1}-x_1,\;\dots,\;T_r^{\,n_r}-x_r)
\]
is again regular.
\note{les crochets manquent au tapuscrit et sont rétablis}

In theorem 7.1.\ the normality assumptions as well as reducedness assumptions
are quite superfluous, only lemme 7.2.\ is used (which should be stated in
paragraph 1, maybe in a single proposition with the analoguous statement
concerning normalisty, and properties $(R_k)$ and $(S_k)$). On the other hand
the proof on page 32 seems to me a little short at one point; I tried to
indicate some details in the margin.

As already stated in my notes, the statement of 9.4.\ is not the most natural
one; one should use here the notion of tame ramification in the broad sense
(just one divisor $D$ being given, whith normal crossings relatively to $S$).
\struck{This is} The proof would, I guess, require to use the results of
sections 7 and 8 in a slightly different form (using the suitable variant of
the notion of tame ramification, cf \struck{example} particular case b) in my
notes). This however does not make any difficulty, in fact the proofs given in
7 and 8 work for any kind of ramification data, provided we assume a statement
of the type 7.2.\ in paragraph 7.

In 10.2.\ and 10.3.\ the significance of the result comes from the fact that
the groups considred are just $\pi_1(U,\;)$ resp.\ $\pi_1(U,\;)^{(p)}$, and the
theorems should be stated accordingly.

I hope you will be able to get along with these rather sketchy indications.
Please do not hesitate to call back on me for whatever does not seem clear to
you!

With best wishes,

\section*{« Données de ramification sur un schéma » (pages 12 à 20)}

\page{12}
\textbf{1. Données de ramification sur un schéma.}

Soit $S$ un schéma (préschéma dans l'ancienne terminologie), $\underline{Et}_S$
le site étale de $S$. Nous appellerons \emph{donnée de ramification sur $S$}
\struck{la donnée d}, $\underline{R} = $ un couple $(\underline{C},
\underline{s})$, où $\underline{C}$ est une catégorie fibrée sur
$\underline{Et}_S$ (\struck{en particulier} i.e.\ c'est une catégorie
$\underline{C}$ munie d'un foncteur fibrant $p : \underline{C} \to
\underline{Et}_S$) et où $\underline{s} : \underline{C} \to \underline{Fl}_S$
est un foncteur cartésien, où $\underline{Fl}_S$ désigne la catégorie fibrée
des \struck{faisceaux} schémas relatifs sur les objets $S'$ de
$\underline{Et}_S$. On fera attention qu'on ne suppose pas que la catégorie
fibrée $\underline{C}$ soit un champ (i.e.\ qu'on puisse recoller morphismes et
objects de $\underline{C}$). Un objet $M$ de la catégorie fibre
$\underline{C}_{S'}$ ($S' \in \mathrm{Ob}\,\underline{Et}_S$) sera appelé un
\emph{modèle de ramification} pour $\underline{R}$, et le schéma
$\underline{s}(M)$ sur $S'$ sera appelé la \emph{réalisation géométrique} de ce
modèle. On fera attention qu'on ne suppose pas le foncteur $\underline{s}$
fidèle ni pleinement fidèle, donc les morphismes entre modèles de ramification
ne peuvent être identifiés aux morphismes entre leurs réalisations
géométriques. Dans les exemples envisagés par la suite, les $\underline{s}(M)$
sont finis sur $S'$.
\note{dans tout le tapuscrit les flèches sont laissées en blanc, la machine
n'en portant pas ; elles sont rétablies ici sans autre mention}

Considérons le foncteur composé $\underline{C} \xrightarrow{\ \underline{s}\ }
\underline{Fl}_S \xrightarrow{\ \text{source}\ } (\mathrm{Sch})$, d'où un
foncteur (le foncteur « réalisation géométrique »)
\[
  \underline{C} \longrightarrow (\mathrm{Sch}) .
\]
Considérons d'autre part la catégorie fibrée sur $(\mathrm{Sch})$ des schémas
relatifs étales et finis (dont la catégorie fibre en $T$ est donc la catégorie
des revêtements étales finis de $T$), et soit $\underline{D}$ la catégorie
fibrée sur $\underline{C}$ image inverse de la catégorie fibrée précédente par
le foncteur réalisation géométrique. La fibre de cette catégorie fibrée en
l'objet $M$ de $\underline{C}$ est donc la catégorie des revêtements étales
finis de la réalisation géométrique $\underline{s}(M)$ de $M$. On appellera
\emph{revêtement de}

\page{13}
Appelons \emph{crible couvrant} de $\underline{C}$ toute sous-catégorie
\emph{pleine} $\underline{C}_0$ de $\underline{C}$ qui est un crible (i.e.\
pour toute flèche, l'origine appartient à $\underline{C}_0$ s'il en est de même
de l'extrémité) et tel que les $p(M)$ pour $M \in \mathrm{Ob}\,\underline{C}_0$
recouvrent l'objet final $S$ de $\underline{Et}_S$. Pour tout $\underline{C}_0$,
soit $\mathrm{Rev}^{\underline{R}}_{\underline{C}_0}(S)$ la catégorie des
sections cartésiennes de $\underline{D}$ au dessus de $\underline{C}_0$. Pour
$\underline{C}_0$ variable, ces catégories forment \struck{\ill{}} un système
inductif de catégories \struck{\ill{}} (NB si $\underline{C}'_0 \subset
\underline{C}_0$, on obtient un foncteur restriction évident
$\mathrm{Rev}^{\underline{R}}_{\underline{C}_0}(S) \to
\mathrm{Rev}^{\underline{R}}_{\underline{C}'_0}(S)$), \struck{et} dont la
limite inductive est notée $\mathrm{Rev}^{\underline{R}}(S)$ et appelée
\emph{catégorie des revêtements de $S$ de type $\underline{R}$}.

Signaons que dans tous les cas que nous aurons à utiliser, les foncteurs de
transition $\mathrm{Rev}^{\underline{R}}_{\underline{C}_0}(S) \to
\mathrm{Rev}^{\underline{R}}_{\underline{C}'_0}(S)$ sont pleinement fidèles,
donc aussi les foncteurs $\mathrm{Rev}^{\underline{R}}_{\underline{C}_0}(S)
\to \mathrm{Rev}^{\underline{R}}(S)$, qui permettent donc d'identifier (à
équivalence près) les $\mathrm{Rev}^{\underline{R}}_{\underline{C}_0}(S)$ à des
sous-catégories pleines de la catégorie $\mathrm{Rev}^{\underline{R}}(S)$ des
revêtements de type $S$. Nous supposerons \struck{de plus} que pour tout
$\underline{C}_0$, \struck{le foncteur composé} et tout $X \in
\mathrm{Ob}\,\mathrm{Rev}^{\underline{R}}_{\underline{C}_0}(S)$, le foncteur
composé $\underline{C}_0 \xrightarrow{\ X\ } \underline{D}
\xrightarrow{\ \text{source}\ } (\mathrm{Sch})_{/S}$, (où le deuxième foncteur
est le foncteur déduit du foncteur source sur la catégorie des schémas finis
étales relatifs,) admet une limite inductive dans $(\mathrm{Sch})_{/S}$, que
nous noterons $\underline{r}(X)$, et que pour $\underline{C}'_0 \subset
\underline{C}_0$, désignant par $X'$ la restriction de $X$ à
$\underline{C}'_0$, l'homomorphisme canonique $\underline{r}(X') \to
\underline{r}(X)$ est un isomorphisme. (Ces conditions seront satisfaites dans
les exemples 1 à 3 ci-dessous, \emph{de généralité croissante}.) Les différents
foncteurs $\underline{r}$ se raccorderont alors et

\page{14}
type $\underline{R}$ de $S$ une section cartésienne de $\underline{D}$ sur
$\underline{C}$. Ces revêtements de type $\underline{R}$ forment une catégorie
\struck{notée suivant les notations standard $\varinjlim \underline{D}/\underline{C}$)}
que nous noterons $\mathrm{Rev}_{\underline{R}}(S)$. On définira la
\emph{réalisation géométrique} d'un revêtement $\Phi$ de type $\underline{R}$
comme étant la limite inductive (si elle existe) dans $(\mathrm{Sch})$ du
foncteur composé $\underline{C} \xrightarrow{\ \Phi\ } \underline{D} \to
(\mathrm{Sch})_{/S}$, où le deuxième de ces foncteurs est le foncteur source
\struck{associant à tout revêtement} (déduit du foncteur source sur la
catégorie des schémas finis étales relatifs). Si pour tout $\Phi$, cette limite
inductive existe, on obtient un foncteur réalisation géométrique
\[
  \underline{r} : \mathrm{Rev}_{\underline{R}}(S) \longrightarrow (\mathrm{Sch})_{/S} .
\]
\note{$\Phi$ est frappé à la machine comme un o barré, faute de capitale
grecque ; le sens est celui d'un revêtement de type $\underline{R}$}
\note{plusieurs traits à l'encre bleue portent la phrase « Si pour tout $\Phi$,
cette limite inductive existe, » plus haut dans le paragraphe ; la
transposition n'est pas exécutée}

Dans les cas qui nous intéresseront par la suite, ce foncteur $\underline{r}$
sera fidèle, mais pas nécessairement pleinement fidèle; il importera donc de
bien distinguer entre une revêtement de type $\underline{R}$ de $S$, et sa
réalisation géométrique, le premier n'étant pas déterminée par la seule
connaissance du second.

Avant d'aller plus loin, donnons quelques exemples.

\emph{Exemple 1.} Donnons-nous un schéma relatif $Z$ sur $S$, muni d'un groupe
fini $G$ opérant (\emph{à droite disons}) par $S$-automorphismes. Soit
$\underline{C}_0$ la catégorie réduite à un objet et dont le monoïde des
endomorphismes est $G$, et prenons pour $\underline{C}$ la catégorie produit
$\underline{Et}_S \times \underline{C}_0$, \struck{pour foncteur $\underline{s}$
le foncteur $(S',g)$ \ill{}}. La donnée d'un foncteur cartésien de
$\underline{C}$ dans une catégorie fibrée $\underline{C}'$ sur
$\underline{Et}_S$ équivaut manifestement à la donnée d'un objet de la
catégorie fibre $\underline{C}'_S$ avec une opération de $G$ dessus, de sorte
que $(Z,G)$ nous fournit bien un foncteur cartésien $\underline{s}$. On trouve
aussitôt que la catégorie $\mathrm{Rev}_{\underline{R}}(S)$ est équivalente
canoniquement à la catégorie des revêtements finis étales $Z'$ de $Z$, munis
d'une opération (droite) de $G$ \struck{\ill{}} rendant le morphisme structural
$Z' \to Z$ un $G$-morphisme.

\page{15}
La réalisation géométrique d'un tel $Z'$ existe si et seulement si le quotient
$Z'/G$ existe, \emph{et elle s'identifie alors à $X = Z'/G$.} Il en est ainsi
si toute orbite de $G$ dans $Z$ est contenue dans un ouvert affine, ce qui
implique en effet qu'il en est de même de toute orbite de $G$ dans $Z'$.
\struck{Intuitivement, la considération \ill{} les revêtements de type
$\underline{R}$ correspondent à des revêtements} Dans les cas que nous avons en
vue, $Z$ est un revêtement de $S$, donc la condition sur les orbites sera
automatiquement satisfaite (il suffit que $Z$ soit affine sur $S$)
\struck{\ill{}} de plus les réalisations géométriques $X = Z'/G$ seront
automatiquement finies sur $S$, \struck{donc} \emph{i.e.} des revêtements. En
première approximation, ce seront les revêtements de $S$ qui (en \struck{un}
un sens convenable \dots) deviennent étales au dessus du revêtement (ramifié en
général) donné $Z$, ce qui donne un sens intuitif à la notion de « donnée de
ramification ».

Soit maintenant $U$ un ouvert de $S$ tel que $Z|U$ soit un revêtement principal
de groupe $G$, et supposons que $Z|U = V$ soit dense dans $Z$ (plus
généralement, que toute partie ouverte et fermée \emph{et stable par $G$} de
$Z$ contenant $Z|U$ soit égale à $Z$). Alors le foncteur $Z' \mapsto Z'|V$ de
la catégorie des $G$-revêtements étales de $Z$ dans la catégorie des
$G$-revétments étales de $V$ est fidèle, d'autre part le foncteur $Z'_V \mapsto
Z'_V/G$ de la catégorie des $G$-revêtements étales de $V$ dans la catégorie des
revêtements étales de $U$ est une équivalence de catégories, à fortiori fidèle,
d'où on conclut que le foncteur \emph{réalisation géométrique} $Z' \mapsto
Z'/G = X$ est \emph{fidèle}.

Suivant le yoga général, \struck{\ill{}} \emph{la catégorie des} revêtements de
type $\underline{R}$ est alors équivalente à la catégorie des « revêtements de
$S$ munis d'une structure de type $\underline{R}$ », où par « structure de type
$\underline{R}$ » sur un revêtement $X$ de $S$ on entend la donnée d'une classe
(mod isomorphisme) d'un $Z' \in \mathrm{Rev}_{\underline{R}}(S)$ et d'un
isomorphisme de $\underline{r}(Z')$ ($= Z'/G$) avec $X$. (Cette

\page{16}
dernière remarque aurait dû figurer avant l'exemple, elle s'applique chaque
fois que le foncteur $\underline{r}$ est fidèle). \struck{Nous pouvons également}
Supposons maintenant $V$ dense dans $Z$ et de plus $Z$ \emph{normal}, alors on
voit aussitôt que le foncteur $Z' \mapsto Z'|V$ est pleinement fidèle, d'où
résulte facilement qu'il en est de même du foncteur réalisation géométrique
$Z' \mapsto Z'/G$; donc dans ce cas, la catégorie $\mathrm{Rev}_{\underline{R}}(S)$
est équivalente à une sous-catégorie pleine de la catégorie des revêtements
normaux de $S$, \emph{étales au dessus de $U$, et} dont toute composante
irréductible domine une composante irréductible de $S$ rencontrant $U$.

\emph{Exemple 2.} C'est une simple variante de l'exemple précédent : au lieu de
se donner un $S$-schéma $Z$ à groupe fini d'opérateurs $G$, on se donne un
système projectif $(Z_i, G_i)_{i \in I}$ de tels objets ($I$ ensemble ordonné
filtrant). \struck{La catégorie $\underline{R}$} Le système projectif des
groupes $(G_i)$ définit une catégorie fibrée $\underline{C}_0$ sur la catégorie
associée à $I$ (les fibres de $\underline{C}_0$ étant les catégories à un seul
objet définies par les $G_i^{\,0}$), et on prendra $\underline{C} =
\underline{Et}_S \times \underline{C}_0$, et on définit $\underline{s}$ de façon
évidente. La catégorie $\mathrm{Rev}_{\underline{R}}(S)$ est maintenant
équivalente à la catégorie $\varinjlim$ des catégories
$\mathrm{Rev}_{\underline{R}_i}(S)$ définies séparément par les $(Z_i, G_i)$.

Un cas intéressant est celui où les foncteurs transition (qui s'interprètent de
façon évidente comme des foncteurs images inverses de revêtements étales à
opérateurs) sont pleinement fidèles, de sorte que les
$\mathrm{Rev}_{\underline{R}_i}(S)$ s'identifient (à équivalence près) à des
sous-catégories pleines de $\mathrm{Rev}_{\underline{R}}(S)$ qui épuisent
évidemment cette dernière. Il en sera ainsi si les morphismes de transition
$(Z_i, G_i) \to (Z_j, G_j)$ sont tels que $G_i \to G_j$ soit surjectif et
$Z_i \to Z_j$ induit un isomorphisme $Z_j \simeq Z_i/N_{ij}$, où $N_{ij}$ est
le noyau de $G_i \to G_j$ (en effet, il en résulte que pour tout revêtement
étale $Z'_j$

\page{17}
\marginal{$G' \to G$ \quad $M' \to M$}
\note{les deux flèches, à l'encre, sont au verso de la feuille ; elles fixent
le sens des morphismes de l'exemple 3 du recto, où $G' = \mathrm{Aut}(M')$ va
dans $G = \mathrm{Aut}(M)$ au-dessus de $M' \to M$}

\page{18}
de $Z_j$, on a $Z'_j \simeq Z'_{ij}/N_{ij}$, où $Z'_i$ est l'image inverse de
$Z'_j$ sur $Z_i$). Dans ce cas, le foncteur « réalisation géométrique » est
bien défini, et induit surchaque $\mathrm{Rev}^{\underline{R}_i}(S)$ le
foncteur réalisation géométrique déjà explicité dans l'exemple 1.
\struck{Bien sûr, les conditions} Notons que les critères signalés dans
l'exemple 1 pour que le foncteur réalisation géométrique soit fidèle resp.\
pleinement fidèle, s'étendent de façon immédiate au cas envisagé ici.

\emph{Exemple 3.} Supposons que $\underline{R} = (\underline{C}, \underline{s})$
satisfasse aux conditions suivantes. Le foncteur $\underline{s}$ se factorise
en $\underline{C} \xrightarrow{\ \underline{s}'\ } \underline{Flop}_S \to
\underline{Fl}_S$, où $\underline{Flop}_S$ désigne la catégorie fibrée sur
$\underline{Et}_S$ des schémas relatifs à groupe fini (opérant à droite)
d'opérateurs, et $\underline{Flop}_S \to \underline{Fl}_S$ est le foncteur
« oubli des opérations ». \struck{Pour tout \ill{} $M$ \ill{} $\underline{s}$
\ill{}} \emph{Il faut préciser que (en vue de l'exemple kummérien plus bas)}
nous prenons comme morphismes d'un schéma à opérateurs $(Z,G)$ sur $S$ dans un
autre $(Z',G')$ les couples $(f,u)$, où $f : X \to X'$ est un $S$-morphisme et
$u : G_S \to G'_S$ un morphisme (localement constant, \emph{pas nécessairement
constant !}) de schémas en groupes sur $S$, de telle façon que $f$ soit
compatible avec $u$ dans le sens évident.
\note{$f : X \to X'$ : le tapuscrit écrit $X$, $X'$ là où les objets sont $Z$,
$Z'$ ; la faute est sur la page et n'est pas corrigée ici}

On suppose que pour toute flèche $M' \to M$ dans un $\underline{C}_S$,
\struck{si} la flèche correspondante $\underline{s}(M') = (Z',G') \to
\underline{s}(M) = (Z,G)$ induit un isomorphisme $Z \simeq Z'/H$, où $H$ est le
noyau de $G'_S \to G_S$ (qui est un sous-groupe localement constant de
$G'_S$). On suppose \struck{\ill{}} aussi que pour tout objet $M$ de
$\mathrm{Ob}\,\underline{C}_{S'}$, posant $\underline{s}(M) = (Z,G)$, on se
donne un \struck{automorphisme} homomorphisme $G \to \mathrm{Aut}(M)$,
\struck{et pour tout $g \in G$, il existe un endomorphisme $g$ de $M$ uniquement
déterminé}, tel que pour tout $g \in G$, si $\bar{g}$ est son image,
\struck{minée, tel que} $\underline{s}(\bar{g})$ soit le couple
$(g_Z, \mathrm{int}(g))$, et que tout endomorphisme de $M$ soit localement de
la forme $\bar{g}$; on \struck{veut} suppose de plus que
\struck{la} l'homomorphisme $G \to \mathrm{Aut}(M)$ soit compatible avec images
inverses, et que pour toute flèche $M' \to M$ dans $\underline{C}_{S'}$,
considérons $M'$ (resp.\ $M$) comme muni ud groupe $G'$ ($G$) d'opérateurs,
$M' \to M$ soit compatible avec $G' \to G$.

\page{19}
\note{la feuille reprend l'exemple 3 sous une autre forme ; les trois premières
lignes, frappées à part en tête, remplacent la fin de la page précédente. Les
deux moutures sont conservées telles quelles}

Dans ce cas, le foncteur « réalisation géométrique » est bien défini et induit
sur chaque $\mathrm{Rev}_{\underline{R}_i}(S)$ le foncteur réalisation
géométrique déjà explicité dans l'exemple précédent.

de $Z_j$, on a $Z'_j \simeq Z'_i/N_{ij}$, où $Z'_i$ est l'image inverse de
$Z'_j$ sur $Z_i$). Notons que les critères signalés dans l'exemple 1 pour que
le foncteur $\underline{r} : \mathrm{Rev}_{\underline{R}}(S) \to
(\mathrm{Sch})_{/S}$ soit fidèle resp.\ pleinement fidèle, s'étendent de façon
immédiate au cas envisagé ici.

\emph{Exemple 3.} Supposons que $\underline{R} = (\underline{C}, \underline{s})$
satisfasse aux conditions suivantes : pour tout $S' \in
\mathrm{Ob}\,\underline{Et}_S$, et $M \in \mathrm{Ob}\,\underline{C}_{S'}$, le
groupe des automorphismes $G$ de $M$ est fini, et pour toute flèche $M' \to M$
de $\underline{C}_{S'}$, il existe un unique morphisme de groupe $G' =
\mathrm{Aut}(M')$ dans $G = \mathrm{Aut}(M)$ tel que $M' \to M$ soit compatible
avec les opérations de $G'$, de plus si $H$ est son noyau, le morphisme
$\underline{s}(M') \to \underline{s}(M)$ induit un isomorphisme
$\underline{s}(M')/H \simeq \underline{s}(M)$. \struck{Enfin} On suppose
\struck{que} de plus que pour deux morphismes de $M'$ dans $M$,
\struck{\ill{}} les \struck{\ill{}} l'un est déduit de l'autre par composition
avec un élément de $G$. Enfin, on suppose qu'étant donnés deux objets $M$, $M'$
de $\underline{C}_{S'}$, il existe une famille surjective de morphismes étales
$S'_i \to S'$ telle que pour tout $i$, $M_{S'_i}$ et $M'_{S'_i}$ soient majorés
par un objet $M''_i$ de $\underline{C}_{S'_i}$.

Utilisant la théorie de la descente, on voit sous ces conditions que le
foncteur réalisation géométrique $\mathrm{Rev}_{\underline{R}}(S) \to
(\mathrm{Sch})_{/S}$ est défini, et à valeurs dans les \struck{schémas finis
sur} revêtements de $S$ (en supposant que les $\underline{s}(M)$ sont des
revêtements, bien sûr; pour pouvoir appliquer la descente, il suffirait même
que les $\underline{s}(M)$ soient affines sur $S'$). Ici encore, on a des
critères pour que ce foncteur soit fidèle resp.\ pleinement fidèle : il suffit
que pour tout $S' \in \mathrm{Ob}\,\underline{Et}_S$ et tout $M \in
\mathrm{Ob}\,\underline{C}_{S'}$, la réalisation géométrique
$\underline{s}(M)$, munie du groupe d'automorphismes $G = \mathrm{Aut}(M)$,
satisfasse aux conditions \emph{correspondantes} envisagées dans l'exemple 1
pour $(Z,G)$.

\emph{Cas particuliers du type kummérien} : a) Supposons données des sections

\page{20}
\struck{$(a_i)_{i \in I}$ de \ill{} sur $\underline{S}$} une famille
$(D_i)_{i \in I}$ de diviseurs (de Cartier) sur $S$. Pour tout $S'$ étale sur
$S$, on désigne par $\underline{C}_{S'}$ la catégorie suivante : les objets
sont les \struck{couples} \emph{quadruples $M = $} $(\underline{a},
\underline{n}, G, \underline{u})$, où $\underline{a} = (a_i)_{i \in I}$ est une
famille d'équations pour les $D_{iS'}$, $\underline{n} = (n_i)_{i \in I}$ une
famille d'entiers $\geqslant 1$ premiers aux caractéristiques résiduelles de
$S'$, \struck{\ill{}} $G$ un groupe fini, et $u : G_{S'} \xrightarrow{\ \sim\ }
\uncertain{\mu}_{\underline{n}}\,S'$ un isomorphisme de schémas en groupes sur
$S'$. Étant donné un deuxième tel quadruple $M' = (\underline{a}',
\underline{n}', G', u')$, \struck{un morphisme de second \ill{} le premier est
$M$ dans $M$} $\mathrm{Hom}_{S'}(M', M)$ est vide sáuf si $\underline{n}' \geqslant
\underline{n}$, auquel cas c'est l'ensemble des $\underline{b} = (b_i)_{i \in I}$
tels que $\underline{b}^{\,\underline{n}} = \underline{a}'\underline{a}^{-1}$.
\note{« 2tant donné » sur la page, pour « Étant donné » : la machine ne porte
pas de capitale accentuée. La lettre grecque de $\mu_{\underline{n}}$ est
laissée en blanc et restituée d'après le contexte kummérien}

La composition des morphismes, et la définition des foncteurs changements de
base pour une flèche $S'' \to S'$ dans $\underline{Et}_S$, sont clairs. Le
foncteur cartésien $\underline{C} \to (\mathrm{Sch})_{/S}$ est défini en
associant à $M = (\underline{a}, \underline{n}, G, \underline{u})$ le
revêtement Kummérien élémentaire $Z^{\underline{n}}_{\underline{a}}$,
\struck{et \ill{}} considéré comme un revêtement à groupe d'opérateurs $G$,
grâce à l'isomorphisme $G \simeq \mu_{\underline{n}}$. On définit de façon
évidente la valeur de $\underline{r}$ sur les flèches de
$\underline{C}_{S'}$, et sa compatibilité avec les changements de base. On voit
alors qu'on est sous les conditions de l'exemple 3. De plus, le foncteur
$\underline{r}$ est fidèle, et il est pleinement fidèle lorsque $S$ est
\struck{noethérien aux points de} localement noethérien, et de plus normal en
les points des $\mathrm{supp}\,D_i$, et si de plus les $D_i$ sont réduits, et
toute composante irréductible d'un $D_i$ est distincte de toute comp.\ d'un
$D_j$ pour $j \neq i$.

b) Supposons donné un diviseur $D$ sur $S$, \struck{\ill{}} et considérons le
faisceau sur $\underline{Et}_S$ associé au préfaisceau qui, à tout $S'$,
associe l'ensemble des \struck{diviseurs} ensembles $\theta$ \emph{$(D_i)$} de
diviseurs sur $S'$ ayant les propriétés suivantes : $D_{S'} = \ill{}\, D_i$, et
pour tout $s' \in S'$, les restrictions à \uncertain{$\mathrm{Spec}$}$(\underline{O}_{S',s'})$
des $D_i$ qui sont non nulles, sont deux à deux disjointes. Soit donnée une
section du

\end{document}
